\chapter{Dialogo Testuale con le
Tradizioni}\label{dialogo-testuale-con-le-tradizioni}

\begin{quote}
\textbf{Scopo dell'apêndice.} Espandere la mappa concettuale del Cap.
13.3 in dieci dialoghi brevi e diretti con le tradizioni che la DEQ$^2$
eredita, modifica o supera. Ciascun dialogo è strutturato come
\emph{obiezione attesa} dalla tradizione + \emph{risposta DEQ$^2$} +
\emph{concessione esplicita} di ciò che la DEQ$^2$ \emph{non} fa com relação a
quella tradizione. La forma è quella del dialogo immaginario: non
pretende di parlare per gli autori citati, pretende di \emph{farsi
parlare} dalle loro categorie.
\end{quote}



\section{\texorpdfstring{Con Giovanni Arrighi (\emph{Il lungo XX
secolo},
1994)}{Con Giovanni Arrighi (Il lungo XX secolo, 1994)}}\label{con-giovanni-arrighi-il-lungo-xx-secolo-1994}

\textbf{Arrighi:} ``Lei ha individuato una transizione 2028-2031 senza
riconoscere che è la \emph{fase autunnale} del ciclo finanziario USA ---
già descritta nel mio ciclo hegemônico
genovese-olandese-britannico-americano. Non sta scoprendo nulla, sta
misurando ciò che io ho narrato.''

\textbf{DEQ$^2$:} ``La concessione è piena: Lei ha narrato la \emph{forma}
della transizione; io provo a \emph{misurarne il bordo}. La
\(T_s^{(2),\text{sys}} \to 0\) del Cap.~7 sull'USA è la traduzione
formale del Suo \emph{signal crisis} (anni '70-'80) seguito da
\emph{terminal crisis} (anni 2020-2030). Quello che aggiungo: la
transizione \emph{non genera automaticamente} un nuovo egemone --- può
generare em vez disso un'asimmetria multi-attore in cui \(\Phi\) esportabile
risiede in un Bilanciere (Brasil) anziché in un nuovo egemone (China).
Lei pensava em termos de \emph{successione}; io penso em termos de
\emph{bifurcação}.''

\textbf{Concessione.} Senza Arrighi, il Cap.~7 e il Cap.~13.7 sarebbero
illeggibili. La DEQ$^2$ \emph{eredita} l'horizonte temporal e l'hipótese di
transizione; \emph{aggiunge} la quantificazione e la possibilità di un
esito non-egemonico.



\section{\texorpdfstring{Con Peter Turchin (\emph{Secular Cycles}, 2009;
\emph{End Times},
2023)}{Con Peter Turchin (Secular Cycles, 2009; End Times, 2023)}}\label{con-peter-turchin-secular-cycles-2009-end-times-2023}

\textbf{Turchin:} ``La Sua \(\Omega\) è la mia \emph{elite
overproduction}; il Suo \(\varepsilon_{\text{dis}}\) è la mia
\emph{Popular Immiseration}. Ha riformulato in notazione greca quello
che ho già modellato cliodinamicamente.''

\textbf{DEQ$^2$:} ``Ha ragione su due variáveis su quattro. Ma manca
\(\Phi\). Lei modella il \emph{denominatore} della stabilità sistemica
(involução + immiserazione \(\Rightarrow\) instabilità) ma non il
\emph{numeratore} (coerência de fase \(\times\) antifragilidade
\(\Rightarrow\) stabilità positiva). Senza \(\Phi\), la cliodinamica
predice \emph{quando} un sistema crolla, ma non \emph{cosa lo riempie}.
La differenza è cruciale per il 2031: due sistemi possono avere la
stessa \(\Omega + \varepsilon_{\text{dis}}\) e finire in M1 (decoerência
con \(\Phi\)-engine) o in M2 (decoerência degenere). La cliodinamica è
agnostica fra M1 e M2; la DEQ$^2$ no.''

\textbf{Concessione.} Le serie temporali storiche di Turchin sono
essenziali per calibrare \(A\) via stress-test retrospettivo (Ap.~C.3).
La DEQ$^2$ rinuncia a fare previsioni a tre secoli (orizzonte cliodinamico)
--- si limita a un \emph{ciclo} (2026-2031).



\section{\texorpdfstring{Con Immanuel Wallerstein (\emph{The Modern
World-System},
1974-2011)}{Con Immanuel Wallerstein (The Modern World-System, 1974-2011)}}\label{con-immanuel-wallerstein-the-modern-world-system-1974-2011}

\textbf{Wallerstein:} ``Il Brasil è semi-periferia ascendente del
sistema-mondo capitalista. Lei lo presenta come \(\Phi\)-engine --- isto é,
come un attore con prerogative \emph{centrali} di produzione di ordine.
È un'illusione: la semi-periferia \emph{non} genera coerenza globale, è
generata da essa.''

\textbf{DEQ$^2$:} ``Concedo che lo schema centro-periferia ha potere
descrittivo notevole per il periodo 1500-2000. Ma la sua linearità
(centro \(\to\) semi-periferia \(\to\) periferia) è un \emph{vincolo}
che la DEQ$^2$ rifiuta di importare. La \(\Phi\) non è correlata alla
posizione nel sistema-mondo: dipende dalla struttura interna del sistema
sociale di un paese (Cap.~4.5, Cap.~5.2). Il Brasil può essere
\emph{semi-periferico per output} e \emph{centrale per \(\Phi\)
esportabile} --- è esattamente la posizione estrutural del Cap.~9. Lei
vincola; io disvincolo.''

\textbf{Concessione.} Lo schema centro-periferia resta utile per
descrivere il \emph{flusso materiale} (capitale, manodopera, risorse).
Ma il \emph{flusso di coerenza} segue regole diverse --- l'errore
estrutural del wallersteinismo è confondere i due flussi.



\section{\texorpdfstring{Con Amado Cervo (\emph{Inserção internacional},
2008)}{Con Amado Cervo (Inserção internacional, 2008)}}\label{con-amado-cervo-inseruxe7uxe3o-internacional-2008}

\textbf{Cervo:} ``L'Estado Logístico è una categoria che ho costruito
per descrivere il Brasil post-2003. Lei lo importa, ma la matematica
DEQ$^2$ rischia di trasformare una \emph{categoria politica vivente} in un
\emph{parâmetro statico}. La politica estera brasileira non è un
coefficiente.''

\textbf{DEQ$^2$:} ``Concessione totale. La
\(\Pi^{\text{Logístico}} = \sum_k M_e^{(k)} Q_0^{(k)} (1 - \delta_R^{(k)} z_R^{(k)})\)
del Cap.~9.1 \emph{non sostituisce} la Sua categoria --- la
\emph{traduce} in linguaggio operativo per la dashboard SPC-DEQ. La
traduzione è utile \emph{operativamente} (permette monitoraggio in tempo
reale) ma è \emph{parziale concettualmente} (perde la dimensione storica
accumulativa che Lei evidenzia). Per questo l'Estado Logístico nella
DEQ$^2$ è \emph{condicional} alla continuità politica --- esattamente la
fragilidade del Cap.~9.5 + Cap.~12.5.''

\textbf{Concessione.} La DEQ$^2$ rinuncia a teorizzare il \emph{processo
storico} di formazione dell'Estado Logístico (1990-2003). Si limita a
operativizzarne i parâmetros \emph{date le condizioni del 2025-2031}.



\section{\texorpdfstring{Con Nassim Nicholas Taleb (\emph{Antifragile},
2012)}{Con Nassim Nicholas Taleb (Antifragile, 2012)}}\label{con-nassim-nicholas-taleb-antifrágil-2012}

\textbf{Taleb:} ``Lei prende la mia antifragilidade individuale e la
incolla a un sistema egemonico globale. La mia teoria è
\emph{anti-sistemica}: l'antifragilidade \emph{vive nei dettagli e nelle
code}, non nelle medie aggregate. Il Suo \(\langle A\rangle\) è
esattamente l'errore che combatto.''

\textbf{DEQ$^2$:} ``Concessione parziale. Sì, \(\langle A\rangle\) è una
media --- ma è modulata da \(\Phi\), e l'asimmetria
\(\Phi\)-moltiplica-\(A\) vs \(\Phi\)-divide-\(\Omega\) è esattamente la
firma matematica della Sua intuizione: il sistema \emph{amplifica} la
fragilidade più di quanto amplifichi l'antifragilidade. La DEQ$^2$ formalizza
la convessità che Lei ha portato come euristica. Inoltre il Cap.~12.6
(\(S_{\mathcal{R}_-}\)) è interamente \emph{taleb-iano}: opzionalità,
ridondanza, disacoplamento dalla traiettoria sistemica. La DEQ$^2$
accoglie la Sua critica trasformandola in protocollo.''

\textbf{Concessione.} La DEQ$^2$ \emph{non} è una teoria della
\emph{previsione} nel senso strong (rifiutato da Taleb); è una teoria
dello \emph{spazio degli scenari} con probabilità soggettive esplicite.
Sotto questo profilo è coerente con \emph{Black Swan} e \emph{Skin in
the Game}.



\section{\texorpdfstring{Con Antonio Gramsci (\emph{Quaderni del
carcere}, 1929-1935) e Robert Cox (\emph{Production, Power, and World
Order},
1987)}{Con Antonio Gramsci (Quaderni del carcere, 1929-1935) e Robert Cox (Production, Power, and World Order, 1987)}}\label{con-antonio-gramsci-quaderni-del-carcere-1929-1935-e-robert-cox-production-power-and-world-order-1987}

\textbf{Gramsci/Cox:} ``L'egemonia non è \(T_s^{(2),\text{sys}}\). È
rapporto fra forze materiali, istituzionali e culturali in cui il
consenso e la coercizione si combinano \emph{organicamente}. Lei riduce
l'organico al numerico.''

\textbf{DEQ$^2$:} ``Concedo che \(T_s^{(2),\text{sys}}\) è una
\emph{misura} dell'egemonia, non la \emph{spiegazione} dell'egemonia. La
spiegazione resta gramsciana: il consenso (\(s_{iii}\) DEQ$^2$) e la
coercizione (\(\sigma \cdot \varepsilon_{\text{dis}}\)) si combinano in
proporzioni che variano per blocco storico. Quello che la DEQ$^2$ aggiunge
è la \emph{condizione di stabilità}: per qualunque blocco storico,
\(T_s^{(2),\text{sys}} > 1\) è condizione necessaria di durabilità. La
DEQ$^2$ è portanto \emph{interna} alla teoria gramsciana dell'egemonia, non
alternativa.''

\textbf{Concessione.} La nozione di \emph{blocco storico} gramsciano
contiene una temporalità \emph{qualitativa} (cesura, conservazione,
rivoluzione passiva, restaurazione) che la DEQ$^2$ non cattura. La
transição de fase del Cap.~3.5.3 corrisponde solo alla \emph{cesura}
gramsciana; le altre forme richiedono estensione.



\section{\texorpdfstring{Con Edward Luttwak (\emph{Strategy: The Logic
of War and Peace},
1987-2001)}{Con Edward Luttwak (Strategy: The Logic of War and Peace, 1987-2001)}}\label{con-edward-luttwak-strategy-the-logic-of-war-and-peace-1987-2001}

\textbf{Luttwak:} ``La grande strategia è \emph{paradossale}: ciò che
funziona si esaurisce, ciò che non funziona viene ripreso. La Sua
matematica è \emph{lineare} nel tempo. Si perde il paradosso.''

\textbf{DEQ$^2$:} ``Concessione. Le forme \(T_s^{(2),\text{sys}}\) sono
monotoniche (Cap.~3.5.1); la dinamica del Cap.~5.4 ammette però
retroazioni non lineari (R6 \(\Omega \to \varepsilon_{\text{dis}}\), R3
\(\Phi \to \varepsilon_{\text{dis}}\) negativa) che producono
comportamento \emph{paradossale} nell'evoluzione di
\(\boldsymbol{p}(t)\): la stessa azione può essere antifrágil in M1 e
fragile in M2. È il paradosso luttwakiano nella forma DEQ$^2$: la strategia
ottima dipende dallo scenario realizzato, e lo scenario realizzato non è
noto al momento della scelta.''

\textbf{Concessione.} La DEQ$^2$ \emph{non} analizza il \emph{registro
militare} della grande strategia (manovra, deception, logistica). Si
limita al \emph{registro estrutural} (coerenza, fragilidade,
involução). Per la grande strategia operativa Lei resta referência.''



\section{Con la Escola de Frankfurt (Marcuse, Adorno, Habermas) ---
via
Topografia}\label{con-la-scuola-di-francoforte-marcuse-adorno-habermas-via-topografia}

\textbf{Marcuse/Adorno/Habermas:} ``L'apparato critico francofortese ha
denunciato la \emph{razionalità strumentale} come distruttrice della
razionalità sostanziale. Lei produce un'altra razionalità strumentale:
la matematizzazione del campo discorsivo. La denuncia, applicata a Lei
stesso, è devastante.''

\textbf{DEQ$^2$:} ``Concedo l'auto-implicazione. Difesa: la formalizzazione
del campo \(\Pi\) topografico (eredità Topografia \S{}5) \emph{non} riduce
\(z\) --- pelo contrário, \emph{misura quanto \(z\) è presente}. La
\(T_s^{(2),\text{sys}}\) premia sistemi con \(z\) alto (asse della
complexidade dialética al numeratore di \(Q_0\)). Una razionalità
strumental que misura --- e premia --- la razionalità sostanziale è
\emph{forse} l'unica via per parlare alla decisione politica nel 2026
senza essere tagliati fuori dal discorso pubblico. Habermas avrebbe
odiato la DEQ$^2$; ma Habermas pubblicava libri come
\emph{Legitimationsprobleme} destinati ad accademici, non a policymaker.
Il Bilanciere parla alla decisione.''

\textbf{Concessione.} La DEQ$^2$ \emph{non} è critica francofortese --- è
\emph{sintesi tecnica della critica francofortese}. Perde
l'incandescenza dialettica originaria; guadagna in azionabilità.
Trade-off accettato.



\section{\texorpdfstring{Con Xiang Biao (\emph{Self as Method}, 2020;
\emph{Anti-involution},
2021-2024)}{Con Xiang Biao (Self as Method, 2020; Anti-involution, 2021-2024)}}\label{con-xiang-biao-self-as-method-2020-anti-involution-2021-2024}

\textbf{Xiang Biao:} ``Il \emph{Neijuan} è una categoria etnografica
viva, nata nei dormitori delle università cinesi negli anni 2010 e
diffusa nei contesti 996. Lei la trasforma in \(\Omega\) --- un
parâmetro di equação. Le persone che vivono il Neijuan non vivono
\(\Omega\); vivono fatica.''

\textbf{DEQ$^2$:} ``Concessione totale. Il passaggio dal Neijuan
etnografato a \(\Omega\) misurato è una \emph{perdita di contenuto
fenomenologico}. Difesa: la perdita è funzionale a un obiettivo diverso
dal Suo. La Sua etnografia \emph{fa sentire} il problema; la DEQ$^2$
\emph{rende contabile} il problema. Le due operazioni sono
complementari, non sostitutive. Inoltre il Cap.~8.2 cita Lei
\emph{direttamente} come \emph{firma empirica} della variável --- non
come decorazione retorica. Senza \emph{Self as Method}, il Cap.~8
sarebbe un esercizio aggregato senza ancoraggio empirico.''

\textbf{Concessione.} La DEQ$^2$ \emph{non} può sostituire la ricerca
etnografica. Le serve come \emph{lettore di senso}; senza, leggerebbe
deflazione PPI senza vedere persone.''



\section{\texorpdfstring{Con Yoshiki Kuramoto (\emph{Chemical
Oscillations, Waves, and Turbulence},
1984)}{Con Yoshiki Kuramoto (Chemical Oscillations, Waves, and Turbulence, 1984)}}\label{con-yoshiki-kuramoto-chemical-oscillations-waves-and-turbulence-1984}

\textbf{Kuramoto:} ``Il mio modello descrive oscillatori chimico-fisici.
Voi sociologi lo applicate da decenni a sistemi sociali con disinvoltura
preoccupante. Le hipótese (campo médio, simmetria della distribuzione
\(g(\omega)\), acoplamento sinusoidale) \emph{non} sono verificaçãobili
sui sistemi sociali.''

\textbf{DEQ$^2$:} ``Concessione: l'Apêndice A.2 dichiara esplicitamente A1
(Kuramoto applicabile) come \emph{hipótese operativa}, non come fatto
verificaçãodo. Difesa: la matematica di Kuramoto è la \emph{più semplice}
descrizione di sincronização fra oscillatori non identici. Per \(N\)
grande e accoppiamenti deboli, qualunque modello di sincronização
converge in qualche limite a Kuramoto (universalità). La DEQ$^2$ è portanto
\emph{minimalista}: usa la più povera matematica di sincronização
disponibile, e ne dichiara i limiti. Estensioni a Kuramoto multi-cluster
(Strogatz 2003) sono pianificate per il paper EN --- proprio per i
fenomeni che il modello base non cattura, come la polarizzazione USA.''

\textbf{Concessione.} La DEQ$^2$ \emph{non} dimostra che Kuramoto è la
matematica corretta del sociale. Argomenta che è una \emph{primeiro
approssimazione consistente}. La verificação empirica è demandata al paper
EN --- coerente con la posizione popperiana del Cap.~11.''



\section{Sintesi dei Dieci Dialoghi}\label{sintesi-dei-dieci-dialoghi}

\Proved{} \textbf{Pattern emergente.} In ciascuno dei dieci dialoghi:

\begin{itemize}
\tightlist
\item
  la DEQ$^2$ fa una \textbf{concessione esplicita} (cosa \emph{non} fa
  com relação alla tradizione);
\item
  la DEQ$^2$ fa una \textbf{risposta sostanziale} (cosa \emph{aggiunge}
  alla tradizione);
\item
  la DEQ$^2$ fa una \textbf{clausola di apertura} (estensioni / verifiche
  pianificate per il futuro).
\end{itemize}

\Hypoth{} \textbf{Posizione metodologica.} La DEQ$^2$ non è teoria
\emph{contro}; è teoria \emph{con}. Questo è coerente con l'hipótese di
campo médio: ogni teoria è un \emph{oscillatore} nel campo
intellettuale; la coerência de fase del campo
\(\Phi^{\text{intellettuale}}\) è promossa da chi accoglie e modifica,
non da chi denuncia e sostituisce. La forma stessa di questa apêndice è
\emph{applicazione riflessiva} della DEQ$^2$ alla propria posizione nel
campo del sapere.

\Proved{} \textbf{Conseguenza editoriale.} I lettori specialisti di una
qualunque delle dieci tradizioni elencate trovano nella DEQ$^2$: 1.
\emph{Riconoscimento esplicito} della loro tradizione (cosa accolto); 2.
\emph{Argomentazione esplicita} del valore aggiunto (cosa nuovo); 3.
\emph{Limite esplicito} del valore aggiunto (cosa lasciato a loro).



\section{Tradizioni Non Dialogate}\label{tradizioni-non-dialogate}

\Conj{} \textbf{Esclusioni esplicite.} Per ragioni di spazio, la DEQ$^2$
\emph{non} dialoga in questa apêndice con:

\begin{itemize}
\tightlist
\item
  \textbf{Foucault} (potere capillare, biopolitica) --- il dialogo
  richiederebbe una riformulazione del campo \(\Pi\) em termos de
  \emph{dispositivo} anziché di \emph{acoplamento}. Rinviato a futura
  estensione.
\item
  \textbf{Latour} (Actor-Network Theory) --- la DEQ$^2$ assume attori
  discreti; ANT li rifiuta. Trade-off metodologico irrisolto.
\item
  \textbf{Bourdieu} (campo, habitus, capitale) --- il \emph{campo}
  bourdieusiano è strutturalmente affine al campo \(\Pi\) topografico,
  ma con \emph{abito disposizionale} anziché \emph{fase di Kuramoto}.
  Riformulazione possibile, non tentata qui.
\item
  \textbf{Tilly} (contentious politics) --- la dinamica dei movimenti
  sociali è esterna al perimetro DEQ$^2$; integrazione possibile via
  estensione \(\dot\Phi\) con termine di forzatura mobilitazionale.
\end{itemize}

\Proved{} \textbf{Posizione.} Le quattro esclusioni sopra non sono
dimissioni; sono \emph{aperture pianificate} per estensioni successive
del programma DEQ$^2$.



\section{§D.4 --- As Ra\'{i}zes Evolutivas: Hamilton, Wilson-Sober, Nowak, Kropotkin}\label{d4-radici-evolutive}

\begin{quote}
\textbf{Prop\'{o}sito desta se\c{c}\~{a}o.} Os quatro di\'{a}logos que
se seguem trazem ao DEQ$^2$ a sua \emph{microfunda\c{c}\~{a}o causal
evolutiva}: a explica\c{c}\~{a}o de \emph{por que} $\Phi$ emerge, de
\emph{como} se erode e de \emph{o que significa} a decoer\^{e}ncia na
escala dos organismos e das popula\c{c}\~{o}es. N\~{a}o \'{e} um
acr\'{e}scimo decorativo --- \'{e} o fundamento causal da estrutura
formal. A deriva\c{c}\~{a}o completa do marco CPGG encontra-se em
\emph{Captured Cooperation and Phase Decoherence} (Marcelino 2026,
working paper).
\end{quote}

\subsection{Com William D. Hamilton (\emph{The Genetical Evolution of
Social Behaviour}, 1964)}\label{con-hamilton-1964}

\textbf{Hamilton:} ``O senhor usa o par\^{a}metro de ordem de Kuramoto
para descrever a coopera\c{c}\~{a}o sist\^{e}mica. Mas de onde vem a
coopera\c{c}\~{a}o? O senhor a pressup\~{o}e como fato; eu a derivo da
gen\'{e}tica. Sem a condi\c{c}\~{a}o $rB > C$, $\Phi > 0$ \'{e} uma
quantidade \'{o}rf\~{a} de causa.''

\textbf{DEQ$^2$:} ``\'{E} a cr\'{i}tica mais profunda que se podia fazer
ao marco original --- e est\'{a} correta. A resposta reside na
\emph{correspond\^{e}ncia estrutural} (n\~{a}o identidade) entre suas
condi\c{c}\~{o}es gen\'{e}ticas e as condi\c{c}\~{o}es sist\^{e}micas.
A condi\c{c}\~{a}o $rB > C$ se traduz no sistema geopol\'{i}tico como:
a coopera\c{c}\~{a}o ($\Phi > 0$) \'{e} est\'{a}vel quando o benef\'{i}cio
coletivo da sincronia supera o custo individual de cooperar --- e isso
tornou-se, no DEQ$^2$ atualizado (Cap.~3, §3.2), a justifica\c{c}\~{a}o
causal de $\Phi$, n\~{a}o apenas sua descri\c{c}\~{a}o. O senhor,
Hamilton, forneceu o \emph{por que} que me faltava.''

\textbf{Concess\~{a}o.} A condi\c{c}\~{a}o $rB > C$ \'{e} aplicada aqui
como \emph{isomorfismo estrutural} (\Hypoth{}) a popula\c{c}\~{o}es de
atores geopol\'{i}ticos, n\~{a}o a popula\c{c}\~{o}es de organismos
biol\'{o}gicos. A transfer\^{e}ncia requer a hip\'{o}tese H1 (§4.0) de
que o campo $\Pi_j$ se mapeia sobre uma fase --- uma hip\'{o}tese
declarada, n\~{a}o um resultado derivado. A deriva\c{c}\~{a}o de Hamilton
\'{e} preservada com plena honestidade epist\^{e}mica.

\subsection{Com David Sloan Wilson e Elliott Sober (\emph{Unto Others},
1994)}\label{con-wilson-sober-1994}

\textbf{Wilson \& Sober:} ``A literatura tratou por muito tempo a
sele\c{c}\~{a}o de grupo como irrelevante, seguindo Dawkins. O senhor
citou inicialmente Dawkins em apoio \`{a} sua mudan\c{c}a de n\'{i}vel
descritivo --- um erro que t\^{e}m visto muitas vezes. A \emph{sele\c{c}\~{a}o
de grupo} n\~{a}o \'{e} uma extens\~{a}o opcional da teoria evolutiva:
\'{e} a estrutura que justifica \emph{qualquer} afirma\c{c}\~{a}o sobre
benef\'{i}cios coletivos.''

\textbf{DEQ$^2$:} ``T\^{e}m inteira raz\~{a}o, e o primeiro rascunho do
marco era defeituoso exatamente neste ponto: Dawkins (1976) rejeita a
sele\c{c}\~{a}o de grupo --- cit\'{a}-lo para justificar a mudan\c{c}a de
n\'{i}vel micro/macro era contradit\'{o}rio. A corre\c{c}\~{a}o est\'{a}
implementada (Cap.~3, §3.2): o fundamento formal \'{e} a teoria da
sele\c{c}\~{a}o multinível, formalizada por Okasha (2006), combinada com
a equa\c{c}\~{a}o de Price (1970) como identidade alg\'{e}brica exata. A
equa\c{c}\~{a}o de Price decompõe $\Delta\bar{z}$ num termo inter-grupo
(proporcional a $\Phi$) e num termo intra-grupo (proporcional a $-\Omega$)
--- uma estrutura que justifica formalmente a forma do numerador e
denominador de $T_s^{(2),\text{sys}}$.''

\textbf{Concess\~{a}o.} A sele\c{c}\~{a}o multinível DEQ$^2$ opera sobre
sistemas socio-institucionais, n\~{a}o sobre alelos. A aplica\c{c}\~{a}o
requer que o ``gen\'{o}tipo'' seja interpretado como configura\c{c}\~{a}o
institucional (\Hypoth{}), e a ``aptid\~{a}o'' como payoff sist\^{e}mico.
Esta reinterpreta\c{c}\~{a}o \'{e} declarada como hip\'{o}tese operativa,
n\~{a}o como fato biol\'{o}gico.

\subsection{Com Martin A. Nowak (\emph{Five Rules for the Evolution of
Cooperation}, 2006)}\label{con-nowak-2006}

\textbf{Nowak:} ``O senhor tem um problema de identifica\c{c}\~{a}o. Seus
atores --- EUA, China, Brasil, Musk --- apresentam simultaneamente
caracter\'{i}sticas de todos os cinco mecanismos cooperativos:
reciprocidade direta (via tratados bilaterais), reciprocidade indireta
(via reputa\c{c}\~{a}o), sele\c{c}\~{a}o de rede (via posicionamento
geopol\'{i}tico), sele\c{c}\~{a}o de grupo (via alian\c{c}as) e
sele\c{c}\~{a}o de parentesco (via afinidade cultural). Como distingue
qual mecanismo guia a din\^{a}mica?''

\textbf{DEQ$^2$:} ``Sua taxonomia dos cinco mecanismos torna-se no
DEQ$^2$ uma \emph{taxonomia diagn\'{o}stica dos atores}, n\~{a}o uma
classifica\c{c}\~{a}o exclusiva dos mecanismos. O CPGG (§D.4, working
paper) foca num mecanismo espec\'{i}fico --- a \emph{captura institucional}
que degrada os cinco mecanismos simultaneamente. A tese \'{e}: quando
$\phi$ (a fra\c{c}\~{a}o extrativa) supera $\phi^*$, os cinco mecanismos
de Nowak ficam simultaneamente enfraquecidos --- a condi\c{c}\~{a}o
$b/c > k_{\text{eff}}(\phi)$ deixa de valer para cada um deles. O DEQ$^2$
n\~{a}o escolhe qual mecanismo opera; diagnostica quando todos
colapsam.''

\textbf{Concess\~{a}o.} O CPGG modela a captura institucional, n\~{a}o a
din\^{a}mica evolutiva de cada um dos cinco mecanismos separadamente. Uma
modelagem completa exigiria cinco CPGGs separados ou um modelo unificado
de intera\c{c}\~{a}o entre mecanismos. Isto \'{e} F1 e F2 na agenda
formal aberta.

\subsection{Com Pyotr Kropotkin (\emph{Mutual Aid: A Factor of
Evolution}, 1902)}\label{con-kropotkin-1902}

\textbf{Kropotkin:} ``Eu documentei, esp\'{e}cie ap\'{o}s esp\'{e}cie e
cultura ap\'{o}s cultura, que a coopera\c{c}\~{a}o \'{e} a norma e a
competi\c{c}\~{a}o a exce\c{c}\~{a}o. O senhor constr\'{o}i um formalismo
para medir a decoer\^{e}ncia --- que \'{e} exatamente o que chamei de
\emph{ruptura da ajuda m\'{u}tua}. Mas vinte anos depois, o Neijuan e a
captura extrativa parecem esmagadores. O senhor esqueceu minha tese
central?''

\textbf{DEQ$^2$:} ``N\~{a}o, eu a \emph{refundei}. Sua intui\c{c}\~{a}o
--- $\Phi > 0$ como estado natural, a decoer\^{e}ncia como estado produzido
--- est\'{a} agora formalizada na Proposi\c{c}\~{a}o K1 (§8 do marco
CPGG): na aus\^{e}ncia de captura institucional ($\sigma_{\text{san}} =
\sigma_0$, $\kappa_E = 0$), a fra\c{c}\~{a}o extrativa de equil\'{i}brio
natural $\phi_{\text{nat}}$ satisfaz $\phi_{\text{nat}} < \phi^*(\sigma_0)$
para par\^{a}metros biologicamente plaus\'{i}veis. A coopera\c{c}\~{a}o
\'{e} o resultado est\'{a}vel; a decoer\^{e}ncia requer a condi\c{c}\~{a}o
produzida $\alpha\phi + \beta\kappa_E > \log(\sigma_0/\sigma_{\text{crit}})$.
O senhor, Kropotkin, estava errado sobre a inevitabilidade da
coopera\c{c}\~{a}o na hist\'{o}ria --- n\~{a}o porque a coopera\c{c}\~{a}o
seja rara, mas porque a captura \'{e} mais f\'{a}cil do que a abertura.
Sua \emph{ajuda m\'{u}tua} \'{e} a condi\c{c}\~{a}o padr\~{a}o, n\~{a}o
o resultado garantido.''

\textbf{Concess\~{a}o.} A invers\~{a}o epist\^{e}mica --- decoer\^{e}ncia
como estado produzido, coopera\c{c}\~{a}o como estado natural --- corre o
risco de tornar-se normativamente ing\^{e}nua: ``baseada na natureza,
portanto devida''. A \emph{fal\'{a}cia naturalista} \'{e} gerida fundando
a parte normativa do Cap.~12 nos princ\'{i}pios de design de Ostrom (1990)
em vez de em Kropotkin diretamente. Ostrom fornece um fundamento
emp\'{i}rico: os commons \emph{podem} ser geridos cooperativamente, sob
condi\c{c}\~{o}es institucionais precisas --- n\~{a}o \emph{naturalmente}
ou \emph{inevitavelmente}. Kropotkin \'{e} inspira\c{c}\~{a}o
fenomenol\'{o}gica; Ostrom \'{e} fundamento normativo.



\section{Sintesi dell'Apêndice}\label{sintesi-dellapêndice}

\Proved{} \textbf{Risultato.} La DEQ$^2$ si posiziona esplicitamente in
quattordici tradizioni intellettuali contemporanee (dieci originali +
quattro evolutive in §D.4), dichiarando per ciascuna concessione, risposta
e clausola di apertura. La forma del dialogo immaginario è
\emph{applicazione riflessiva} dell'apparato (campo médio Kuramoto
applicato al campo intellettuale). Quattro tradizioni ulteriori (Foucault,
Latour, Bourdieu, Tilly) sono dichiarate come estensioni pianificate.

\Hypoth{} \textbf{Tesi conclusiva del libro.} La DEQ$^2$ esiste \emph{fra}
tradizioni, non \emph{contro} di esse --- coerentemente con la sua tesi
sostanziale che la \(\Phi\) vera si genera per acoplamento spontaneo
di sotto-attori, non per imposizione gerarchica. Il libro stesso è un
\emph{atto \(\Phi\)-engine} nel campo intellettuale: produce coerenza
ammettendo dissonanza, e si offre alla critica come condizione di vita.



\emph{Apêndice di posizionamento. Stile: dialogo immaginario,
concessioni esplicite, aperture dichiarate.}
